\documentclass[../../../include/open-logic-section]{subfiles}

\begin{document}

\olfileid{sth}{z}{milestone}
\olsection{$\Z^-$: एक महत्त्वाचा टप्पा}

पुढील विभागात \stagesinf{} कडे पुन्हा वळू. पण अनंततेचे स्वयंसिद्धक मिळाल्यावर आपण एक महत्त्वाचा टप्पा गाठला आहे. आता $\Zminus$ या उपपत्तीसाठी लागणारी सर्व स्वयंसिद्धके आपल्याकडे आहेत. नेमके सांगायचे तर:

\begin{defn}
$\Zminus$ या उपपत्तीची स्वयंसिद्धके अशी आहेत: विस्तारात्मकता, संयोग, जोड्या, घातसंच, अनंतता आणि विभक्तीकरण-रूपबंधाची सर्व उदाहरणे.
\end{defn}

या नावातील \emph{झर्मेलो} म्हणजे झर्मेलो यांची संच उपपत्ती; \emph{वजा} केलेली गोष्ट पुढे पाहू. झर्मेलो यांना हे श्रेय उचित आहे, कारण त्यांनी \citeyear{Zermelo1908Untersuchungen} मधील कामात ही उपपत्ती मूलत: मांडली होती.\footnote{इतिहास आणि तांत्रिक तपशीलांवरील रंजक टिप्पण्यांसाठी \citet[Appendix
A]{Potter2004} पाहा.}

या उपपत्तीत गणिताचा फार मोठा भाग मांडता येतो. विशेषत: \olref[sfr][][]{part} पुन्हा पाहून, तेथे अनौपचारिकरीत्या केलेल्या प्रत्येक गोष्टीची अधिक औपचारिक मांडणी $\Zminus$ मध्ये करता येते, याची तुम्ही \emph{खात्री करून घ्यायला हवी}. (थोडे असे केल्यावर पुढील \olref[sth][z][arbintersections]{sec} कडे थेट वळावेसे वाटू शकते.) म्हणून पुढे वेगळी नोंद न करता आपण किमान $\Zminus$ मध्ये काम करत आहोत, असे मानू.

\end{document}
